\documentclass[11pt,letterpaper]{article}
\usepackage[T1]{fontenc}
\usepackage{lmodern,microtype}
\usepackage[letterpaper,textwidth=6.15in,top=0.96in,bottom=0.98in,headheight=15pt,headsep=18pt]{geometry}
\usepackage{amsmath,amssymb,amsthm,mathtools,booktabs}
\usepackage[table]{xcolor}
\definecolor{inkblue}{HTML}{194D60}
\usepackage{fancyhdr,needspace}
\usepackage[colorlinks=true,linkcolor=inkblue,citecolor=inkblue,urlcolor=inkblue]{hyperref}
\usepackage{bookmark}
\usepackage{xurl}
\hypersetup{pdftitle={Jacobi truncation and the motion of a counted zero pair},pdfauthor={Draft prepared for Edward Baker},pdfsubject={Integer exterior moments, physical-time forcing and signed gap estimates in the Newman heat flow}}
\numberwithin{equation}{section}
\newtheorem{theorem}{Theorem}[section]
\newtheorem{proposition}[theorem]{Proposition}
\newtheorem{lemma}[theorem]{Lemma}
\newtheorem{corollary}[theorem]{Corollary}
\theoremstyle{definition}\newtheorem{definition}[theorem]{Definition}
\theoremstyle{remark}\newtheorem{remark}[theorem]{Remark}
\pagestyle{fancy}\fancyhf{}
\fancyhead[L]{\footnotesize Jacobi truncation and a counted zero pair}
\fancyhead[R]{\footnotesize Edward Baker}
\fancyfoot[C]{\small\thepage}
\renewcommand{\headrulewidth}{0.25pt}
\setlength{\emergencystretch}{1.3em}
\allowdisplaybreaks[2]
\raggedbottom
\begin{document}
\thispagestyle{plain}
\begin{center}
{\small\scshape Analytical working draft\par}
\vspace{9pt}
{\LARGE\bfseries Jacobi truncation and the motion\\[3pt] of a counted zero pair\par}
\vspace{7pt}
{\large Integer exterior moments and signed gap estimates\\ in the Newman heat flow\par}
\vspace{10pt}
{\large Edward Baker\par}
\vspace{2pt}
{\small 20 November 2026\par}
\end{center}
\begin{abstract}
A finite Jacobi matrix gives a Gaussian quadrature for the theta measure in the Newman heat flow. As the measure evolves, the first omitted Jacobi link determines the heat defect of this quadrature. We follow that defect to a pair of Fourier zeros isolated by a contour. It acts on every integer exterior moment through the variation of the field and the motion of the two poles of the moment kernel. The resulting equations remain regular at coalescence in symmetric coordinates. A resolvent difference then gives the exact change of the finite field when two sites are added. For the corresponding change of squared half-gap, we isolate a signed linear contour integral and bound its full nonlinear remainder, including the displacement of the pair center. An explicit compact positive density yields increments of both signs infinitely often on each degree parity. For the canonical theta source, a separate interval certificate resolves a negative degree-24 correction on one physical-time interval. The sharper direct enclosure on that same interval identifies the remaining question of how much complete-source information a finite degree step can replace.
\end{abstract}
\vspace{0.35em}
\setcounter{tocdepth}{1}
{\small\tableofcontents}
\clearpage

\section{The approximation problem}\label{sec:intro}
Consider two zeros of an entire function inside a disk centered on the real axis that contains no other zeros. Their center and their squared half-difference are analytic functions of the coefficients as long as the boundary stays clear. This remains true when the zeros meet. For a field real on the real line, the squared half-difference is positive for a distinct real pair, zero for a double root, and negative for a conjugate pair off the line. A signed estimate for the error in this one quantity is therefore more informative than a bound on the distances traveled by the individual zeros.

In the Newman heat flow, the pair also has an exact equation of motion. Its complete exterior moments were developed in \cite{cm}, and a direct signed contour comparison on an earlier cluster appears in \cite{gt}. If \(D\) denotes its squared half-gap and \(M\) the first exterior moment of all the other zeros, the parent investigation gives
\begin{equation}
 \dot D=2-4MD . \label{eq:intro-clock}
\end{equation}
The constant is the contribution from the two zeros to one another. The product \(MD\) accounts for the rest of the divisor. Equation~\eqref{eq:intro-clock} and the contour formula for \(M\) are equations (2.5)--(2.7) of Baker~\cite{agc}. In that work the approximation is made on the arithmetic side of the heat flow. Here we approximate its positive source measure by Gaussian quadrature. This choice leaves the Fourier observation fixed and allows us to vary physical time at fixed degree and to add sites at fixed time.

Two formulas describe the link between the source and the pair. If \(G_N\) is the finite Fourier transform and \(\beta_j\) the squared Jacobi links, its heat defect is \(\beta_N\beta_{N-1}\partial_{\beta_{N-1}}G_N\). If a target field has relative error \(q\) against this head on a circle enclosing the pair, the linear change of \(D\) is the contour moment of \(-(z-c)q\), with \(c\) the reference pair center. Only the accuracy of that signed moment enters to first order. Theorem~\ref{th:signed} gives the constants in the quadratic remainder and includes the case where the enclosed zeros coincide.

The two operations answer different questions. Varying time exposes a boundary term in the finite Jacobi chain. Increasing the degree reveals an exactly computable correction from two finite chains. Neither operation alone determines the difference from the complete field. We follow the boundary term through the integer hierarchy in Section~\ref{sec:hierarchy}, compute the finite degree difference in Section~\ref{sec:degree}, and convert relative field information into a signed gap estimate in Section~\ref{sec:signed}. The compact-source theorem of Section~\ref{sec:parity} explains why positivity of the quadrature weights cannot supply the missing sign. Sections~\ref{sec:time} and~\ref{sec:evidence} show how analytic time bounds turn the estimate into a uniform statement on one canonical chart. Section~\ref{sec:open} records the outstanding challenges.

\section{The moving Jacobi approximation}\label{sec:source}
\subsection{One source and two evolutions}
With the physical time convention of Polymath~\cite{polymath} and
\cite[(1.1)--(1.5)]{agc}, let
\begin{align}
 \Phi(u)&=\sum_{n\ge1}(2\pi^2 n^4 e^{9u}-3\pi n^2 e^{5u})
                       e^{-\pi n^2 e^{4u}},\qquad u\ge0, \notag\\
 H(t,z)&=\int_0^\infty e^{tu^2}\Phi(u)\cos(zu)\,du,\qquad
 F(t,z)=H(t,z)/H(t,0). \label{eq:source}
\end{align}
Extend \(\Phi\) by reflection. Each summand is positive on the half-line and the full source decays doubly exponentially. Differentiation under the integral gives \(H_t=-H_{zz}\). At \(t=0\), \(8H(0,z)=\xi((1+iz)/2)\) with the usual completed Riemann xi function. For real \(t\),
\[
 d\mu_t(u)=\frac{e^{tu^2}\Phi(|u|)}{2H(t,0)}\,du,\qquad
 F(t,z)=\int_{\mathbb R}e^{izu}\,d\mu_t(u).
\]
The cosine and exponential forms agree by parity. Write \(m_k(t)=\int u^k\,d\mu_t\). Differentiation of the normalization yields
\begin{equation}
 \partial_t d\mu_t=(u^2-m_2)d\mu_t,\qquad
 (\partial_t+\partial_z^2+m_2)F=0 . \label{eq:norm-heat}
\end{equation}

Let \(P_j(u;t)\) be monic orthogonal polynomials and \(h_j=\int P_j^2\,d\mu_t\), so \(h_0=1\). Parity cancels the diagonal coefficient of the three-term recurrence. With \(P_{-1}=0\), \(P_0=1\), it reads
\begin{equation}
 uP_j=P_{j+1}+\beta_jP_{j-1},\qquad
 \beta_0=0,\quad \beta_j=h_j/h_{j-1}\ (j\ge1).
 \label{eq:monic}
\end{equation}
In particular, \(\beta_1=m_2\) and \(h_j=h_j^*:=\prod_{k=1}^j\beta_k\).
In the orthonormal basis \(p_j=P_j/\sqrt{h_j}\), multiplication by \(u\) has off-diagonal entries \(a_j=\sqrt{\beta_j}\) between sites \(j-1\) and \(j\). Its first \(N\) rows and columns define the real symmetric matrix \(J_N(t)\). We call it the \emph{head of degree \(N\)} and observe it through
\begin{equation}
 G_N(t,z)=(e^{izJ_N(t)})_{00}=\sum_{\nu=1}^N\omega_\nu(t)e^{izx_\nu(t)} .
 \label{eq:G}
\end{equation}
Here \(x_\nu\) are the eigenvalues, the roots of \(P_N\), and \(\omega_\nu\) are the squared first coordinates of normalized eigenvectors. The vector at site zero is cyclic since all the links are positive. Thus the nodes are simple, the weights are positive and their sum is one. No closed walk from site zero can visit site \(N\) and return in fewer than \(2N\) steps. Therefore \(m_k=(J_N^k)_{00}\) for \(k<2N\), the Gaussian quadrature identity \cite[\S3.5(v)--(vi)]{dlmf}. The nodes are points in the source coordinate \(u\). The zero pairs studied below belong to the entire function \(z\mapsto G_N(t,z)\).

\begin{proposition}[The boundary term in physical time]\label{prop:defect}
The squared links and the moving readout of degree \(N\ge2\) satisfy
\begin{align}
 \dot\beta_j&=\beta_j(\beta_{j+1}-\beta_{j-1}), \label{eq:volterra}\\
 R_N:=(\partial_t+\partial_z^2+\beta_1)G_N
     &=\beta_N\beta_{N-1}\,\partial_{\beta_{N-1}}G_N .
 \label{eq:finite-defect}
\end{align}
In the last derivative the remaining entries of the link vector are held fixed.
\end{proposition}
\begin{proof}
Because \(\partial_t P_j\) has degree less than \(j\), its inner product with \(P_j\) vanishes. Equations~\eqref{eq:norm-heat} and~\eqref{eq:monic} give
\[
 (\log h_j)'=h_j^{-1}\|uP_j\|_{\mu_t}^2-\beta_1
            =\beta_{j+1}+\beta_j-\beta_1.
\]
Subtract consecutive equations to obtain~\eqref{eq:volterra}, the symmetric quadratic-tilt flow \cite[\S2.3]{vanassche}.

To see what truncation does, fix a time \(t_0\) and take the \(N\)-node quadrature measure at that instant. Tilt its own weights by \(e^{(t-t_0)u^2}\) and renormalize. At \(t_0\) its Jacobi links agree with the original links through \(N-1\) by Gaussian exactness. Its next monic polynomial vanishes at every node, so its squared link \(\beta_N\) is zero. The same norm differentiation applies to the positive norms through \(h_{N-1}\). The two velocities therefore differ only in \(\dot\beta_{N-1}\), by \(\beta_{N-1}\beta_N\). The Fourier transform of the tilted finite measure obeys~\eqref{eq:norm-heat}; its value, spatial derivatives and scalar \(\beta_1\) at \(t_0\) are those of \(G_N\). The chain rule in the positive link coordinates gives~\eqref{eq:finite-defect}. This proof also covers \(N=2\).
\end{proof}

The defect can be evaluated without differentiating zero locations. Write \(\mathcal R(\xi)=(\xi-J_N)^{-1}\). Cofactors along the unique chain of off-diagonal entries give
\[
 \mathcal R_{0,N-1}=\frac{a_1\cdots a_{N-1}}{P_N(\xi)},\qquad
 \mathcal R_{0,N-2}=\frac{\xi a_1\cdots a_{N-2}}{P_N(\xi)} .
\]
Resolvent differentiation gives \(\partial_{\beta_{N-1}}\mathcal R_{00}=
\mathcal R_{0,N-2}\mathcal R_{0,N-1}/a_{N-1}\). For a positive contour \(\Sigma\) enclosing the source nodes, functional calculus in~\eqref{eq:finite-defect} gives
\begin{equation}
 R_N(t,z)=\frac{h_N}{2\pi i}\oint_\Sigma
                \frac{\xi e^{iz\xi}}{P_N(\xi)^2}\,d\xi . \label{eq:sr}
\end{equation}
The extra link in \(h_N\) is the first one omitted by the readout itself \cite{vcd}.

\subsection{The complete quadrature error}
A second identity measures the field error at a fixed time and explains the factorial scale that enters the degree analysis. Let \(U\) have law \(P_N(u)^2\,d\mu_t(u)/h_N\). Independently choose weights uniformly on the simplex of \(2N+1\) nonnegative coordinates summing to one. Let \(X_N\) be the corresponding convex combination of \(U\) and two copies of each node \(x_\nu\). Then
\begin{equation}
 F(t,z)-G_N(t,z)=L_N(z)E_N(z),\qquad
 L_N(z)=\frac{(-1)^Nh_Nz^{2N}}{(2N)!},\quad E_N(z)=\mathbb E e^{izX_N}.
 \label{eq:full-phase}
\end{equation}
Indeed, interpolate \(u\mapsto e^{izu}\) with its first derivative at every node by a polynomial of degree at most \(2N-1\). Gaussian exactness cancels the polynomial. The remainder at \(u\) is \(P_N(u)^2\) times the divided difference on the doubled node list and \(u\). Iterating the fundamental theorem of calculus on the divided difference gives the simplex integral of the \(2N\)-th derivative, including coincident nodes by continuity. The simplex volume is \(1/(2N)!\). Integrating with respect to \(d\mu_t\) gives~\eqref{eq:full-phase}. All interchanges follow from the doubly exponential source tail. This identity, including the law of the phase variable, is derived in \cite{vcd}. It holds for the compactly supported measures used in Section~\ref{sec:parity} as well. Its leading factor is finite-link data. Evaluation of its expectation in general still involves the full measure.

\section{The forced integer hierarchy}\label{sec:hierarchy}

\subsection{A counted pair and its exterior traces}

Let \(Y\) be \(F\) or a moving head \(G_N\) with \(N\ge2\), and set
\(R=(\partial_t+\partial_z^2+\beta_1)Y\).
On a fixed positively oriented circle \(\Gamma:|z-a|=r\) with
\(a\in\mathbb R\), assume a cleared collar and a count of two with
multiplicity. On compact portions of the chart take a buffer \(b>0\)
so the two occurrences lie in \(|z-a|\le r-b\). We use
\(\langle g\rangle=(2\pi i)^{-1}\oint_\Gamma g(z)\,dz\) and write
\begin{equation}\label{h:coordinates}
\begin{gathered}
 B=Y_z/Y,\qquad c=\tfrac12\langle zB\rangle,\qquad
 D=\tfrac12\langle z^2B\rangle-c^2,\qquad
 w=z-c,\qquad Q=w^2-D,\\
 A=Y/Q,\qquad U=A_z/A,\qquad f=R/Y .
\end{gathered}
\end{equation}
The argument principle makes \(c\) the mean and \(D\) the squared
half-difference of the two zeros. Both are real at real time.
The unit \(A\) extends holomorphically and is zero-free on the disk
and a collar. Each enclosed factor has modulus at least \(b\) on
\(\Gamma\), so \(|Q|\ge b^2\) there, also at \(D=0\). The
coordinates and the traces below are locally analytic in time for
the canonical fields, by the contour formulas and analyticity of
the source and finite moment construction.

For any rational weight \(k\) with poles only at the counted pair and
\(k(z)=O(|z|^{-2})\), put
\begin{equation}\label{h:trace}
 \tau[k]= -\langle kB\rangle=-\langle kU\rangle
        =\sum_{\rho\ \mathrm{exterior}} k(\rho).
\end{equation}
Every exterior zero is included with multiplicity. For the complete
field we use the even product and zero count \(O(x\log(2+x))\)
from \cite[Section~2]{agc}. Conjugating \(J_N\) by the alternating
diagonal sign matrix proves evenness of the head. The positive
spectral weights give \(|G_N(z)|\le e^{K_N|\operatorname{Im}z|}\), where
\(K_N=\sup_t\|J_N(t)\|\) on the compact time interval of the chart.
Jensen's formula at radius \(2x\) and \(G_N(0)=1\) give
\(n_N(x)\le 2K_Nx/\log 2\).
Hadamard factorization at order at most one pairs opposite factors
as \(1-z^2/\rho^2\). Evenness removes the possible linear exponential
and normalization fixes the constant. These products thus have locally
uniformly convergent logarithmic derivatives on the cleared collar
when summed in opposite pairs. Both zero counts give absolute
inverse-square summability. To prove \eqref{h:trace}, one exterior
pole contributes \(-k(\rho)\) to \(\langle kB\rangle\), by summing
the residues of \(k(z)/(z-\rho)\) over the counted poles and using
its vanishing residue at infinity. The rational target contribution
\(kQ_z/Q\) has all its poles inside \(\Gamma\) and zero total residue.
Any constant regularizer likewise integrates to zero.

The regular traces are
\begin{equation}\label{h:st}
 S_n=\tau[Q^{-n}]\quad(n\ge1),\qquad
 T_n=\tau[wQ^{-n}]\quad(n\ge2),\qquad M=S_1.
\end{equation}
Where \(D\ne0\), take a local square root \(s^2=D\) and put
\(z_\pm=c\pm s\), \(h=2s\), \(X=z-z_+\), and \(Z=z-z_-\).
The split traces, including both axes, are
\begin{equation}\label{h:split}
 k_{p,q}=X^{-p}Z^{-q},\qquad
 P_{p,q}=\tau[k_{p,q}],\qquad
 p,q\in\mathbb Z_{\ge0},\quad p+q\ge2 .
\end{equation}
All these weights have the required decay, as do their spatial and
coordinate derivatives on the respective charts.

\subsection{Motion of the pair and differentiation of a trace}

A dot denotes forward physical time. The normalized equation gives
\(Y_t/Y=-Y_{zz}/Y-\beta_1+f\) on the collar and hence
\begin{equation}\label{h:burgers}
 \dot B=-B_{zz}-2BB_z+f_z .
\end{equation}
Differentiating the counted moments gives
\(\dot c=-\langle Y_t/Y\rangle/2\) and
\(\dot D=-\langle wY_t/Y\rangle\). In \(Y_{zz}/Y\), write the
quadratic contribution as \(2/Q+2(Q_z/Q)U\) and the holomorphic
contribution as \(A_{zz}/A\). Since
\(\langle1/Q\rangle=0\), \(\langle w/Q\rangle=1\), and
\(wQ_z/Q=2+2D/Q\), the pair equations are
\begin{equation}\label{h:pair}
 \begin{gathered}
 \dot c=\langle (Q_z/Q)U\rangle+\gamma,\qquad
 \dot D=2-4DM+\delta,\\
 \gamma=-\tfrac12\langle f\rangle,\qquad
 \delta=-\langle wf\rangle .
 \end{gathered}
\end{equation}
The center term before \(\gamma\) is \(U(z_+)+U(z_-)\), with value
\(2U(c)\) at a double. At a double \(f=(R/A)/w^2\), so
\(\gamma=-\tfrac12(R/A)_z(c)\) and \(\delta=-(R/A)(c)\).
In particular \(\dot D=2\) at a double of the complete field.

\begin{proposition}[Full trace derivative]\label{h:forceprop}
For a kernel \(k(z;c,D)\) whose coordinate derivatives satisfy
the conditions in \eqref{h:trace}, the full derivative is
\begin{equation}\label{h:force}
 \begin{split}
 \frac{d}{dt}\tau[k]
   &=\langle k(B_{zz}+2BB_z)\rangle
               +\dot c\,\tau[k_c]+\dot D\,\tau[k_D]+\langle k_z f\rangle,\\
 \varepsilon[k]
   &=\langle k_z f\rangle+\gamma\tau[k_c]+\delta\tau[k_D]\\
   &=\big\langle\big(k_z-\tfrac12\tau[k_c]-w\tau[k_D]\big)f\big\rangle .
 \end{split}
\end{equation}
Here \(\varepsilon[k]\) is the part contributed by the residual
\(R\). For a holomorphic field variation with first derivative \(V\)
on the counted chart, the differentials of \(c,D\) and the contour
trace are \(\gamma,\delta,\varepsilon[k]\) with \(f\) replaced by \(V/Y\).
\end{proposition}
\begin{proof}
Differentiate \(-\langle kB\rangle\) on the fixed circle and insert
\eqref{h:burgers}. The identity
\(-\langle k f_z\rangle=\langle k_z f\rangle\) follows by
integration of \((kf)_z\). The first variation uses
\(dB=(V/Y)_z\) and the same integration by parts for the counted
moments. The formulas keep both changes of the kernel, even when
\(V/Y\) has a holomorphic part on the disk.
\end{proof}

For a head, \eqref{eq:finite-defect} shows that these residual terms are
\(\beta_N\beta_{N-1}\) times the derivatives in the last-link
direction, with the other links fixed. For \(F\), all three vanish.
The gap gradient \(-\langle wV/Y\rangle\) supplies the signed first
variation in Section~\ref{sec:signed}. The finite change in each
kernel enters the degree comparison in Section~\ref{sec:degree}.

\subsection{Split rows and axes}

When \(D\ne0\) the parts of \(\dot z_\pm\) contributed by \(R\) are
\(v_\pm=\gamma\pm\delta/h\), and
\begin{equation}\label{h:split-gap}
 \dot h=4/h-2hM+2\delta/h .
\end{equation}
For a split weight the residual part of \eqref{h:force} is
\begin{equation}\label{h:split-force}
 \begin{split}
 \varepsilon_{p,q}={}&p\big(v_+P_{p+1,q}-\langle k_{p+1,q}f\rangle\big)\\
                    &+q\big(v_-P_{p,q+1}-\langle k_{p,q+1}f\rangle\big).
 \end{split}
\end{equation}
A summand with vanishing prefactor is omitted. For \(m,n\ge1\) set
\[
 C_{m,n}=\sum_{i=2}^m P_{i,0}P_{m+1-i,n}
         +\sum_{j=1}^{n-1}P_{m,j}P_{0,n+1-j}.
\]
On the axes put
\(C_{m,0}=\sum_{i=2}^{m-1}P_{i,0}P_{m+1-i,0}\) and
\(C_{0,n}=\sum_{j=2}^{n-1}P_{0,j}P_{0,n+1-j}\).
An empty sum is zero.

\begin{theorem}[Forced split recurrence]\label{h:splitthm}
Write \(A_h=2/h-hM\). The mixed rows \((p,q\ge1)\) satisfy
\begin{equation}\label{h:mixed}
\begin{split}
 \dot P_{p,q}={}&p\{C_{p+1,q}-(p+3)P_{p+2,q}
                   -(q+2)P_{p+1,q+1}+A_hP_{p+1,q}\}\\
 &+q\{C_{p,q+1}-(p+2)P_{p+1,q+1}
                   -(q+3)P_{p,q+2}-A_hP_{p,q+1}\}
   +\varepsilon_{p,q}.
\end{split}
\end{equation}
The two axes have \(p\ge2\) and \(q\ge2\), respectively, and satisfy
\begin{equation}\label{h:axes}
 \begin{split}
 \dot P_{p,0}&=p\{C_{p+1,0}-(p+3)P_{p+2,0}-2P_{p+1,1}
                          +(2/h)P_{p+1,0}\}+\varepsilon_{p,0},\\
 \dot P_{0,q}&=q\{C_{0,q+1}-2P_{1,q+1}-(q+3)P_{0,q+2}
                          -(2/h)P_{0,q+1}\}+\varepsilon_{0,q}.
 \end{split}
\end{equation}
\end{theorem}

\begin{proof}
We give a contour derivation of the heat terms of
\cite[Eqs.~(3.5)--(3.7)]{agc} and the moving-head corrections of
\cite{jh02}. Set \(B_Q=Q_z/Q\) and
\(\mathcal I[k]=\langle kU^2\rangle\).
With \(B=B_Q+U\), move the derivatives in the heat term of
\eqref{h:force} onto \(k\) by integration by parts to obtain
\begin{equation}\label{h:heat-term}
 \langle k(B_{zz}+2BB_z)\rangle
     =-\tau[k_{zz}-2B_Qk_z]-\mathcal I[k_z].
\end{equation}
The terms containing only \(B_Q\) are rational, have all their
poles enclosed and decay at infinity with zero residue.

For the quadratic term, integrate the divided difference
\(-[k(u)-k(v)]/(u-v)\) against \(U(u)U(v)\), with \(u\) on
\(\Gamma\) and \(v\) on a smaller circle still enclosing the pair.
Both integrals are normalized by \((2\pi i)^{-1}\).
The \(k(u)\) term vanishes by the inner Cauchy integral. Integrating
the \(k(v)\) term first in \(u\) gives \(\mathcal I[k]\).
For \(k=k_{m,n}\), subtracting each reciprocal power gives
\begin{equation}\label{h:split-dd}
 -\frac{k_{m,n}(u)-k_{m,n}(v)}{u-v}
   =\sum_{i=1}^{m} k_{i,0}(u)k_{m+1-i,n}(v)
    +\sum_{j=1}^{n} k_{m,j}(u)k_{0,n+1-j}(v).
\end{equation}
Products integrate as products of the two linear functionals
\(-\langle kU\rangle\). For a first-power factor these functionals
are \(-U(z_+)\) or \(-U(z_-)\). Thus for mixed \(m,n\),
\(\mathcal I[k_{m,n}]=C_{m,n}-(U(z_+)+U(z_-))P_{m,n}\).
When \(m\ge3,n=0\) its two endpoint terms instead give
\(-2U(z_+)P_{m,0}\). The opposite axis gives
\(-2U(z_-)P_{0,n}\) for \(n\ge3\).

In \eqref{h:heat-term} use
\(\partial_z k_{p,q}=-p k_{p+1,q}-q k_{p,q+1}\).
In the block with weight \(l=k_{p+1,q}\), for example, the
expression in braces is
\(\tau[l_z-2B_Ql]+\mathcal I[l]+\dot z_+\tau[l]
 -\langle lf\rangle\).
The full velocity at a simple target is
\(\dot z_+=2/h+2U(z_+)+v_+\), and the other is
\(\dot z_-=-2/h+2U(z_-)+v_-\).
These follow from \eqref{h:pair} using
\(U(z_+)-U(z_-)=-hM\), itself the residue identity for
\(-\langle U/Q\rangle\). The two endpoint terms therefore give
\(\pm A_h\) in mixed blocks and \(\pm2/h\) on the axes.
The linear terms follow from \(B_Q=1/X+1/Z\).
This proves the formulas without differentiating any exterior
zero, and so also covers exterior multiplicities.
\end{proof}

\subsection{Regular rows through a collision}

\begin{theorem}[Forced symmetric recurrence]\label{h:regularthm}
In the ranges in \eqref{h:st}, the equations on the full counted
chart are
\begin{align}
 \dot S_n={}&2n\sum_{j=1}^n S_jS_{n+1-j}
   +2n\sum_{j=2}^n\big(T_jT_{n+2-j}+DS_jS_{n+2-j}\big)\notag\\
 &-4n(n+2)S_{n+1}-4n(n+3)DS_{n+2}+\varepsilon^S_n ,
 \label{h:S-row}\\
 \dot T_n={}&2(2n-1)\sum_{j=2}^n T_jS_{n+1-j}
    +4nD\sum_{j=2}^n T_jS_{n+2-j}\notag\\
 &-4(n^2+n-1)T_{n+1}-4n(n+3)DT_{n+2}+\varepsilon^T_n .
 \label{h:T-row}
\end{align}
The residual terms are
\begin{equation}\label{h:ST-force}
\begin{split}
 \varepsilon^S_n={}&-2n\langle wQ^{-n-1}f\rangle
                      +2n\gamma T_{n+1}+n\delta S_{n+1},\\
 \varepsilon^T_n={}&\langle ((1-2n)Q^{-n}-2nDQ^{-n-1})f\rangle\\
                  &+\gamma\big((2n-1)S_n+2nDS_{n+1}\big)
                      +n\delta T_{n+1}.
\end{split}
\end{equation}
\end{theorem}

\begin{proof}
Both kernels satisfy \(k_c=-k_z\). The even kernel has
\(k_z=-2nwQ^{-n-1}\), \(k_D=nQ^{-n-1}\), and the odd kernel has
\(k_z=(1-2n)Q^{-n}-2nDQ^{-n-1}\), \(k_D=nwQ^{-n-1}\).
This yields \eqref{h:ST-force}.

To use the nested-circle identity from the split proof, write
\(\alpha=Q(u)^{-1}\), \(\lambda=Q(v)^{-1}\) and
\(L_j=\sum_{i=0}^{j-1}\alpha^{j-1-i}\lambda^i\) with \(L_0=0\).
The negative divided differences of \(Q^{-j}\) and \(wQ^{-j}\)
are respectively
\[
 (w(u)+w(v))\alpha\lambda L_j,\qquad
 \alpha\lambda L_{j-1}+
   (D+w(u)w(v))\alpha\lambda L_j .
\]
These follow from \(Q(u)-Q(v)=(u-v)(w(u)+w(v))\).
For this computation only write \(T_1=-\langle wU/Q\rangle\).
Then
\begin{align*}
 \mathcal I[Q^{-j}]&=2\sum_{i=1}^{j} T_iS_{j+1-i},\\
 \mathcal I[wQ^{-j}]&=\sum_{i=1}^{j-1} S_iS_{j-i}
             +\sum_{i=1}^{j}(T_iT_{j+1-i}+DS_iS_{j+1-i}).
\end{align*}
The center speed before adding \(\gamma\) is \(-2T_1\). In the
even row this cancels the two \(T_1\) endpoints from
\(2n\mathcal I[wQ^{-n-1}]\). In the odd row it cancels
the \(T_1\) terms of
\((2n-1)\mathcal I[Q^{-n}]+2nD\mathcal I[Q^{-n-1}]\).
The \(-4DM\) part of the gap speed cancels the \(DS_1\) endpoint
products in the even row and the \(T_{n+1}S_1\) term in the
\(D\)-weighted odd sum. The remaining linear terms come from
\(-\tau[k_{zz}-2B_Q k_z]+2\tau[k_D]\). For \(Q^{-n}\) the first
operator inside \(\tau\), before the minus sign, has coefficients
\(2n(2n+5)\) on \(Q^{-n-1}\) and \(4n(n+3)D\) on \(Q^{-n-2}\).
For \(wQ^{-n}\) they are \(2(n+2)(2n-1)\) on \(wQ^{-n-1}\)
and \(4n(n+3)D\) on \(wQ^{-n-2}\).
Adding the \(2\tau[k_D]\) part gives exactly
\eqref{h:S-row}--\eqref{h:T-row}. The auxiliary \(T_1\) was a
contour value and cancels from the recurrences. Setting \(f=0\)
recovers \cite[Eqs.~(3.11)--(3.12)]{agc}.
\end{proof}

The lowest rows, which enter the degree and signed gap comparisons, are
\begin{equation}\label{h:low}
 \begin{split}
 \dot M &=2M^2-12S_2-16DS_3+\varepsilon^S_1,\\
 \dot S_2 &=8MS_2+4T_2^2+4DS_2^2-32S_3-40DS_4+\varepsilon^S_2,\\
 \dot T_2 &=(6M+8DS_2)T_2-20T_3-40DT_4+\varepsilon^T_2 .
 \end{split}
\end{equation}
Here \(\varepsilon^S_1=-2\langle wQ^{-2}f\rangle
+2\gamma T_2+\delta S_2\). In particular, comparing \(\dot M\)
requires the change of \(12S_2+16DS_3\) along with the change
of this entire residual term.

For \(D\ne0\), \(S_n=P_{n,n}\) and
\(T_n=(P_{n-1,n}+P_{n,n-1})/2\) for \(n\ge2\).
Conversely, for \(m=\max(p,q)\), the numerator of
\(k_{p,q}=(w-s)^{m-p}(w+s)^{m-q}/Q^m\) has degree at most
\(2m-2\) in \(w\). Reducing by \(w^2=Q+D\) expresses each split
trace in admissible \(S,T\) with coefficients polynomial in \(s\).
Near a double \(t_*\), this gives \(P_{p,q}=a_{p,q}(t)+s(t)b_{p,q}(t)\)
with analytic coefficients. Differentiating the kernel at \(s=0\)
gives
\[
 b_{p,q}(t_*)=(p-q)\sum_{\rho\ \mathrm{exterior}}
      (\rho(t_*)-c(t_*))^{-p-q-1}.
\]
If \(b_{p,q}(t_*)\dot D(t_*)\ne0\), the split derivative has a
term of order \(1/s\). The symmetric rows retain all higher
traces and stay regular at \(D=0\) by the contour buffer.

\section{The two-degree increment and physical-time transport}
\label{sec:degree}

Write \(G_-=G_N(t,z)\) and \(G_+=G_{N+2}(t,z)\), with \(N\ge2\). For any
quantity evaluated in the two fields, let \(\Delta X=X_+-X_-\).
The two-degree numerator \(d_G=G_+-G_-\) is determined by a finite
coefficient prefix at each time. Its comparison with the heat residual
requires the next link. The finite formula and the argument on a common
Fourier zero contour follow the degree and time identities in \cite{jh03}.

\subsection{A shared finite numerator}

\begin{theorem}[Two-degree resolvent and simplex formula]
\label{d:finite-numerator}
Suppress \(t\) and write \(P_j(\xi)=\det(\xi I-J_j)\) and
\(s_j(\xi)=((\xi I-J_j)^{-1})_{00}\). As in
\eqref{eq:full-phase}, \(L_j=h_j(iz)^{2j}/(2j)!\) with
\(h_j=\prod_{k=1}^j\beta_k\).
Then
\begin{equation}
 s_{N+2}(\xi)-s_N(\xi)=
       \frac{h_N\xi}{P_N(\xi)P_{N+2}(\xi)}.
 \label{d:resolvent}
\end{equation}
Let \(x_1,\ldots,x_m\), \(m=2N+2\), be the combined list of spectral
nodes of \(J_N,J_{N+2}\), with every occurrence retained. Set
\(X=\sum_{i=1}^m u_i x_i\) for \(u\) uniformly distributed on the
probability simplex and let \(\chi(z)=\mathbb E e^{izX}\).
For a positively oriented spectral contour \(\Sigma\) enclosing the nodes,
\begin{align}
 d_G(t,z)
 &=\frac{h_N}{2\pi i}\oint_\Sigma
           \frac{\xi e^{iz\xi}}{P_N(\xi)P_{N+2}(\xi)}\,d\xi
 \label{d:contour-numerator}\\
 &=L_N\left(\chi(z)+\frac{z\chi'(z)}{2N+1}\right).
 \label{d:simplex-numerator}
\end{align}
In particular, the parenthesized factor is entire, including at
\(z=0\). For \(y_*\ge|\operatorname{Im}z|\) and
\(K\ge\|J_{N+2}\|\), its deviation from \(1\) satisfies
\begin{equation}
 \left|\chi+\frac{z\chi'}{2N+1}-1\right|
 \le |z|^2\left(\frac12+\frac1{2N+1}\right)
     \sigma^2 e^{y_*K},
 \qquad
 \sigma^2=
 \frac{\operatorname{tr}J_N^2+\operatorname{tr}J_{N+2}^2}
      {(2N+2)(2N+3)} .
 \label{d:simplex-bound}
\end{equation}
\end{theorem}

\begin{proof}
Set \(C_j(\xi)=\xi I-J_j\) and write \(a_j=\sqrt{\beta_j}\).
Appending the last site changes the \(00\) resolvent entry by
\[
 \frac{\beta_j(C_j^{-1})_{0,j-1}^{\,2}}
      {\xi-\beta_j(C_j^{-1})_{j-1,j-1}}.
\]
The endpoint cofactor gives
\((C_j^{-1})_{0,j-1}=\prod_{k=1}^{j-1}a_k/P_j\).
The determinant formula for the Schur complement gives the denominator
\(P_{j+1}/P_j\). Hence \(s_{j+1}-s_j=h_j/(P_jP_{j+1})\).
Adding consecutive differences and expanding the last row and column of
\(P_{N+2}\) gives \(P_{N+2}+\beta_{N+1}P_N=\xi P_{N+1}\), which
proves \eqref{d:resolvent} as a rational identity. Integrating the
scalar resolvents against \(e^{iz\xi}\), or applying Cauchy's formula
in an eigenbasis, gives \eqref{d:contour-numerator}.

For the simplex factor put \(p=m-1\), and integrate in the \(p\) free
coordinates \(u_i\ge0\), \(\sum_{i=1}^{p}u_i\le1\), with
\(u_m=1-\sum_{i=1}^p u_i\). The volume is \(1/p!\). For a multi-index
\(\alpha\) with \(|\alpha|=k\), successive beta integrals give
\(\mathbb E\prod_i u_i^{\alpha_i}=p!\prod_i\alpha_i!/(p+k)!\).
Consequently the multinomial formula gives
\(\mathbb E X^k=k!p!H_k(x)/(p+k)!\), where
\(H_k(x)=\sum_{|\alpha|=k}\prod_i x_i^{\alpha_i}\).
On a large circle the Laurent expansion
\(\prod_i(\xi-x_i)^{-1}=\xi^{-m}\sum_{k\ge0}H_k(x)\xi^{-k}\)
therefore gives the Hermite--Genocchi formula
\[
 \frac1{2\pi i}\oint_\Sigma
      \frac{f(\xi)}{\prod_i(\xi-x_i)}\,d\xi
 =\frac{1}{p!}\mathbb E f^{(p)}(X).
\]
Indeed both sides give \(H_k(x)\) for \(f(\xi)=\xi^{p+k}\), and zero
for monomials of degree less than \(p\). The formula follows for entire
\(f\) by locally uniform Taylor expansion and contour deformation.
This proof also covers repeated nodes. For \(f(\xi)=e^{iz\xi}\), the
right side is \((iz)^p\chi(z)/p!\). Applying \(i^{-1}\partial_z\)
inserts the numerator \(\xi\) and gives \eqref{d:simplex-numerator}.

The node lists have zero sum, so \(\mathbb E X=0\). The simplex
second moments are \((1+\mathbf 1_{i=j})/[m(m+1)]\), giving the stated
variance. Compression and convexity give \(|X|\le K\). Taylor's
integral remainders, using \(\mathbb E X=0\), give
\[
 |\chi-1|\le \tfrac12 |z|^2\mathbb E[X^2e^{y_*|X|}],
 \qquad
 |\chi'|\le |z|\mathbb E[X^2e^{y_*|X|}],
\]
which imply \eqref{d:simplex-bound}.
\end{proof}

The poles of the resolvents in \(\xi\) are nodes of the finite source
quadratures. We count zeros of the entire Fourier readouts on a
separate contour in \(z\). Formulas \eqref{d:contour-numerator}
and \eqref{d:simplex-bound} use links through \(\beta_{N+1}\).
By \eqref{eq:finite-defect}, the two normalized heat residuals use
through \(\beta_{N+2}\).

The full-source error laws \eqref{eq:full-phase} give a second
expression for \(d_G\). Put
\(A_j=\beta_{j+1}z^2/[(2j+1)(2j+2)]\), so \(L_{j+1}=-A_jL_j\).
Subtraction at the two degrees gives
\begin{equation}
 d_G=L_N(E_N-A_NA_{N+1}E_{N+2}).
 \label{d:source-numerator}
\end{equation}
An enclosure of \(E_j-1\) by a disk of radius \(e_j\) puts the
parenthesized factor in the disk with center \(1-A_NA_{N+1}\) and
radius \(e_N+|A_NA_{N+1}|e_{N+2}\). The explicit coefficients in
\eqref{d:source-numerator} use through \(\beta_{N+2}\). Evaluating
each exact \(E_j\) requires the full source. The inherited phase
enclosures use their own strip and time hypotheses
\cite{vcd,jh03,jh04}.

\subsection{A single logarithm for both counted coordinates}

Fix \(\Gamma:|z-a|=r\) with real \(a\), oriented positively. On a
compact time interval \(I\), assume each field has two zeros inside
\(\Gamma\), with multiplicities, and the fields share a zero-free
annular collar. Assume both pairs lie in \(|z-a|\le r-b\), where
\(0<b\le r\). Write \(B_j=G_{j,z}/G_j\), \(w_j=z-c_j\),
\(Q_j=w_j^2-D_j\), \(j\in\{-,+\}\). On \(\Gamma\) the individual
root distances are at least \(b\). Thus
\[
 |Q_j|\ge b^2,\qquad |w_j|\le g:=2r-b,\qquad
 |D_j|\le (r-b)^2.
\]
Equal counts give zero winding for \(G_+/G_-\) on the collar, so it
has a holomorphic logarithm \(\ell\). A choice at one time continues
in \(I\) by integration of the time logarithmic derivative on the
collar. All the contour identities are independent of the constant
branch choice.

\begin{proposition}[Logarithmic coordinate difference]
\label{d:log-coordinates-prop}
With \(\bar c=(c_++c_-)/2\),
\begin{equation}
 \Delta c=-\tfrac12\langle\ell\rangle,\qquad
 \Delta D=-\langle(z-c_-)\ell\rangle-(\Delta c)^2
         =-\langle(z-\bar c)\ell\rangle.
 \label{d:coordinate-log}
\end{equation}
If \(L\ge\sup_\Gamma|\ell|\) for some branch, then
\(|\Delta c|\le rL/2\) and \(|\Delta D|\le rgL\).
The formulas extend through a double zero inside the contour.
\end{proposition}

\begin{proof}
For \(p_{k,j}=\langle(z-a)^kB_j\rangle\), the argument principle
gives \(c_j=a+p_{1,j}/2\), \(D_j=p_{2,j}/2-p_{1,j}^2/4\).
Since \(\Delta B=\ell_z\), integration by parts gives
\(\Delta p_k=-k\langle(z-a)^{k-1}\ell\rangle\). Substitution for
\(k=1,2\) gives \eqref{d:coordinate-log}. The norm of a contour
integral is bounded by \(r\) times the supremum on \(\Gamma\), and
\(|z-\bar c|\le g\). Both formulas use counted moments that are analytic
in time on the cleared chart.
\end{proof}

For example, if \(\sup_\Gamma|d_G/G_-|\le\eta<1\), the series
branch of \(\log(1+d_G/G_-)\) gives \(L=-\log(1-\eta)\).
If the straight interpolation \(G_-+s d_G\), \(0\le s\le1\), has
positive clearance on the circle, its log integral in \(s\) gives
\(L=\sup_\Gamma |d_G|/\inf_{s,\Gamma}|G_-+s d_G|\).
These estimates require the indicated denominator floors.

The exterior trace \(M_j=S_{1,j}=-\langle Q_j^{-1}B_j\rangle\)
also changes with degree. More generally, for a fixed-order trace
\(X_j=-\langle k_j B_j\rangle\) of the pair-dependent kernel \(k_j\),
\begin{equation}
 \Delta X=\langle k_{+,z}\ell\rangle
           -\langle(k_+-k_-)B_-\rangle .
 \label{d:kernel-difference}
\end{equation}
This follows by subtracting the two trace integrals and integrating
\(\ell_z\) by parts. For \(M\), the change of the kernel is given by
\[
 Q_+^{-1}-Q_-^{-1}
   =\frac{(2z-c_--c_+)\Delta c+\Delta D}{Q_+Q_-}.
\]
In particular, if \(a_c,a_D\) bound \(|\Delta c|,|\Delta D|\),
then
\begin{equation}
 |\Delta M|\le
 r\left(\frac{2L}{b^3}
        +\frac{(2g a_c+a_D)\sup_\Gamma|B_-|}{b^4}\right).
 \label{d:delta-m-bound}
\end{equation}
Here \(|(Q_+^{-1})_z|\le 2/b^3\) by differentiating its two
linear factors. Both denominators in the kernel change have their
own floors. Formula \eqref{d:kernel-difference} retains the correlation
with \eqref{d:coordinate-log} before taking absolute values.

\subsection{Residuals and endpoint transport}

Let \(f_j=R_j/G_j\), \(V_j=G_{j,zz}/G_j\), and
\(r_j=-\langle w_jf_j\rangle\), the gap force denoted by \(\delta\)
in \eqref{h:pair}.
The common scalar \(\beta_1\) cancels from the time logarithmic
derivative. The two-field heat identity on the collar is
\begin{align}
 \ell_t+\ell_{zz}+(B_++B_-)\ell_z
  &=\Delta f
   =\frac{\Delta R-f_-d_G}{G_+},
 \label{d:two-field-heat}\\
 \ell_t&=\Delta f-\Delta V
   =\frac{\Delta R-d_{G,zz}-(f_--V_-)d_G}{G_+}.
 \label{d:log-time}
\end{align}
Indeed \(V_j=B_{j,z}+B_j^2\) and
\(G_{j,t}/G_j=f_j-V_j-\beta_1\). In the numerator of
\eqref{d:log-time},
\(\Delta R-d_{G,zz}=(\partial_t+\beta_1)d_G\) vanishes to
order at least \(2N\) at the origin by \eqref{d:simplex-numerator}.
This cancellation concerns the Taylor order at \(z=0\).
The residual difference also includes the changed center.
\begin{equation}
 \Delta r=-\langle w_+\Delta f\rangle+\Delta c\,\langle f_-\rangle.
 \label{d:residual-difference}
\end{equation}
Given a floor \(m_+\le\inf_\Gamma|G_+|\) with \(m_+>0\),
\eqref{d:two-field-heat} gives
\(\|\Delta f\|_\Gamma\le(\|\Delta R\|_\Gamma+
\|f_-\|_\Gamma\|d_G\|_\Gamma)/m_+\), where \(\|\cdot\|_\Gamma\)
denotes the circle supremum. Thus
\(|\Delta r|\le r(g\|\Delta f\|_\Gamma+a_c\|f_-\|_\Gamma)\).

For a differentiable collar test \(k\), integration by parts in
\eqref{d:two-field-heat} gives
\[
 \langle k\,\Delta f\rangle
 =\frac{d}{dt}\langle k\ell\rangle
 -\big\langle\big[k_t-k_{zz}+(k(B_++B_-))_z\big]\ell\big\rangle .
\]
The interior contribution is part of the time integral of the residual.
For the complete gap difference the two center terms combine to give
\begin{equation}
 (\Delta D)'=-4\Delta(DM)+\Delta r
   =-\langle (z-\bar c)\ell_t\rangle
           +\bar c'\langle\ell\rangle .
 \label{d:raw-gap}
\end{equation}
The first equality is the difference of the pair equations
\eqref{h:pair}. The second is differentiation of
\eqref{d:coordinate-log}. Both are regular when either \(D_j=0\).

\begin{theorem}[Complete relative forcing on a time subinterval]
\label{d:transport}
For \(u\le v\) in \(I\), set
\[
 W_j(t;v)=\exp\left(-4\int_t^v M_j(s)\,ds\right)
 \qquad(j\in\{-,+\}, t\in[u,v]).
\]
The traces \(M_j\) are real by conjugation symmetry, so both weights
are positive. The full forcing satisfies
\begin{equation}
 \int_u^v W_-(t;v)\,[\Delta r-4D_+\Delta M](t)\,dt
   =\Delta D(v)-W_-(u;v)\Delta D(u).
 \label{d:weighted-gap}
\end{equation}
For \(u<v\), division by \(\int_u^v W_-(t;v)\,dt>0\) gives the
weighted mean of the full forcing on that subinterval. If
\(a_D(t)\ge|\Delta D(t)|\), the absolute value of the integral in
\eqref{d:weighted-gap} is at most \(a_D(v)+W_-(u;v)a_D(u)\).
\end{theorem}

\begin{proof}
Expand \(\Delta(DM)=M_-\Delta D+D_+\Delta M\), so
\((\Delta D)'+4M_-\Delta D=\Delta r-4D_+\Delta M\).
Since \(\partial_t W_-=4M_-W_-\), the integrand in
\eqref{d:weighted-gap} equals \(\partial_t(W_-\Delta D)\).
Evaluation at the endpoints proves the equality, and positivity
of \(W_-\) gives the stated bound.
\end{proof}

The corresponding formulas with each field's own weight are
\[
 \int_u^v W_j(t;v)(2+r_j(t))\,dt=D_j(v)-W_j(u;v)D_j(u).
\]
Subtracting them yields
\(\Delta D(v)-W_+(u;v)\Delta D(u)
 -(W_+-W_-)(u;v)D_-(u)\).
For \(u<v\), each weight at \(u\) is the exponential of \(-4(v-u)\)
times the respective mean of \(M_j\) on \([u,v]\).
For \(a_M\ge|\Delta M|\), the weight difference satisfies
\[
 |W_+(t;v)-W_-(t;v)|
 \le W_-(t;v)\left(\exp\left(4\int_t^v a_M(s)\,ds\right)-1\right).
\]
There is also a complete endpoint identity for the \(M\) row.
Write \(\mathcal F_j=12S_{2,j}+16D_jS_{3,j}\) and denote its
residual contribution in the equation for \(M_j\) by \(\varepsilon_{M,j}\).
The row \(M_j'=2M_j^2-\mathcal F_j+\varepsilon_{M,j}\) in
\eqref{h:low} gives, for
\(A(t)=\exp(-2\int_u^t(M_++M_-))\),
\[
 \int_u^v A(t)(\Delta\varepsilon_M-\Delta\mathcal F)(t)\,dt
   = A(v)\Delta M(v)-\Delta M(u).
\]
Here \(\Delta\mathcal F=12\Delta S_2+
16(D_+\Delta S_3+S_{3,-}\Delta D)\) retains the full change
in the higher traces and in the moving gap.

The absolute bound in Theorem~\ref{d:transport} applies after
integration of the signed sum. Numerical use still requires count
and clearance on the selected chart, bounds for the endpoint log
functionals and enclosure of the weight integral. For example, an
error \(e_P\) in an evaluated integral of \(M_-\) contributes at
most \(\widetilde W_-(e^{4e_P}-1)\) to the error of the evaluated
weight. A direct contour evaluation furnished with the same
logarithm, endpoint data and weight can use the same cancellation.
For comparison with the complete theta field, the logarithm and time
identities still apply on its common counted chart, with \(R_F=0\).
That comparison requires source-error and clearance estimates for
the complete field. A sign conclusion requires an additional signed
estimate on the selected chart.

\section{A signed first variation of the squared half-gap}\label{sec:signed}
We now compare a reference field \(Y=G_N\) with either the complete field \(Z=F\) or the next head \(Z=G_{N+2}\). The argument applies to any two holomorphic fields with the stated geometry. Fix a positive circle \(\Gamma=\{|z-a|=r\}\), \(a\in\mathbb R\), and use \(\langle g\rangle=(2\pi i)^{-1}\oint_\Gamma g(z)\,dz\). Each field has two zeros in the disk, counted with multiplicity, both in \(|z-a|\le r-b\) for \(0<b\le r\). Write \(c_Y,D_Y\) and \(c_Z,D_Z\) for the two centers and squared half-gaps. Assume both fields are real under complex conjugation, and define on a zero-free collar
\begin{equation}
 q=(Z-Y)/Y,\qquad \|q\|_{\infty}\le\eta<1,\qquad K=-\log(1-\eta).
 \label{eq:q}
\end{equation}
The ratio measures the field error relative to the available boundary size. A smaller absolute field error on a near-zero boundary can give a larger \(q\).

\Needspace{14\baselineskip}
\begin{theorem}[Signed change on a counted circle]\label{th:signed}
Let \(q_0\) be a continuous approximation to \(q\) on \(\Gamma\) with \(\|q-q_0\|_{\infty,\Gamma}\le e\). Set
\[
 J=-\langle(z-c_Y)q_0\rangle,\qquad
 \mathcal A=r(2r-b)\{e+K-\eta\}+\tfrac14 r^2K^2 .
\]
Then
\begin{equation}
 |D_Z-D_Y-J|\le\mathcal A . \label{eq:signed}
\end{equation}
In particular, for \(0\le\eta\le1/4\),
\begin{equation}
 \mathcal A\le r(2r-b)e+
          \left(\frac{16r^2}{9}-\frac{2rb}{3}\right)\eta^2 .
 \label{eq:signed-small}
\end{equation}
The estimate includes a double zero of either field.
\end{theorem}
\begin{proof}
The series for \(\ell=\log(1+q)\) gives a single-valued holomorphic logarithm on the collar. The argument principle gives \(c_Y=\tfrac12\langle zY_z/Y\rangle\) and \(D_Y=\tfrac12\langle z^2Y_z/Y\rangle-c_Y^2\). Subtract the corresponding equations for \(Z\), using \(Z_z/Z-Y_z/Y=\ell_z\), and integrate by parts:
\begin{align}
 \Delta c&=-\tfrac12\langle\ell\rangle,\notag\\[-2pt]
 \Delta D&=-\langle(z-c_Y)\ell\rangle-\tfrac14\langle\ell\rangle^2 .
 \label{eq:center-squared}
\end{align}
The square is the change induced by recentering the quadratic moment. In particular, the first variation is precisely \(-\langle(z-c_Y)q\rangle\). Equation~\eqref{eq:center-squared} is the centered-error identity in \cite[(7.2)]{agc}. On a circle it is also the two-Fourier-coefficient identity of \cite[(8.15)]{agc}. Here we retain the signed first variation even when its input is approximate.

The logarithm series gives \(\|\ell-q\|\le K-\eta\) and \(\|\ell\|\le K\). The pair buffer gives \(|c_Y-a|\le r-b\), hence \(|z-c_Y|\le2r-b\) on \(\Gamma\). Apply \(|\langle g\rangle|\le r\|g\|\) to the two terms of \(\Delta D-J\). This proves~\eqref{eq:signed}. For the quadratic bound use
\(K-\eta\le\eta^2/[2(1-\eta)]\le 2\eta^2/3\) and \(K\le4\eta/3\). All formulas use symmetric moments of the two occurrences.
\end{proof}

For a real sign decision one can take the real part of \(J\), or choose \(q_0\) with real symmetry so that \(J\) is real. If \(\Re J+\mathcal A<0\), the reference gap statistic exceeds the target's. A sufficient condition for a distinct real target pair is instead \(D_Y+\Re J-\mathcal A>0\). These tests answer different mathematical questions. Count and clearance enter both. A boundary ratio below one transfers a reference count by Rouché's theorem. To infer a common buffer from counts one must also localize the roots, for example by counting both on an inner circle. Count equality on the outer boundary alone leaves open how close a root comes to it.

\paragraph{Which source information gives \(J\)?}
If the complete numerator is available, take \(q_0=q\) and \(e=0\). For a cheaper reference to the physical field, \eqref{eq:full-phase} suggests \(q_0=L_N/G_N\): any proved \(|E_N-1|\le\epsilon_N\) on the circle allows the choice \(e=\sup_\Gamma |L_N|\epsilon_N/|G_N|\), with a separate collar bound for \(\eta\). The existing explicit theta phase estimates in \cite{vcd} use \(|\Im z|\le1\) on their physical-time range. A wider circle requires a new bound there. Alternatively use \(q_0=(G_{N+2}-G_N)/G_N\), supplied by Section~\ref{sec:degree}. For a finite-to-finite comparison this is the exact quotient. For a finite-to-complete comparison it still leaves the error
\begin{equation}
 e\ \ge\ \sup_{\Gamma}\frac{|F-G_{N+2}|}{|G_N|}.
 \label{eq:remaining}
\end{equation}
Thus a finite increment acquires one signed part of the comparison while specifying exactly which residual must still be estimated.

\paragraph{An endpoint bound for the time integral.}
Write \(\Delta D=D_Z-D_Y\), \(M_Y\) for the reference exterior moment, and \(r_Y=-\langle(z-c_Y)R_Y/Y\rangle\), with the analogous definition for \(r_Z\). In the notation of~\eqref{d:weighted-gap}, let
\(\mathcal P=r_Z-r_Y-4(M_Z-M_Y)D_Z\) and \(W(u,v)=\exp(-4\int_u^v M_Y)\).
On a common counted time interval, application of the theorem at each endpoint gives
\begin{align}
 &\left|\int_u^v W(s,v)\mathcal P(s)\,ds-
             \{J(v)-W(u,v)J(u)\}\right| \notag\\
 &\hspace{40mm}\le \mathcal A(v)+W(u,v)\mathcal A(u).
 \label{eq:endpoint}
\end{align}
Indeed the exact weighted integral is \(\Delta D(v)-W(u,v)\Delta D(u)\). If an enclosure \(|W-\widehat W|\le\omega\), \(\widehat W\ge0\), is used, the error for \(J(v)-\widehat WJ(u)\) is at most
\(\mathcal A(v)+(\widehat W+\omega)\mathcal A(u)+\omega|J(u)|\).
This bound holds on every subinterval without differentiating the approximate signed input. The complete moment \(M_Y\), including its exterior contributions, is needed for the weight. Direct contour evaluation supplied with the same \(q_0,e,\eta,W\) can use the same bound.

\paragraph{A direct comparison when the roots are simple and real.}
Write \(x_1<x_2\) for the target roots, \(y_1<y_2\) for the reference roots, \(h=(y_2-y_1)/2\) and \(B_0=Z-Y\). For the exact quotient, the residue theorem gives
\begin{equation}
 J=-h\left\{\frac{B_0(y_2)}{Y_z(y_2)}
                  -\frac{B_0(y_1)}{Y_z(y_1)}\right\}. \label{eq:residue}
\end{equation}
The same numerator also allows a direct evaluation. For \(d_i=x_i-y_i\), the mean value theorem gives \(d_i=-B_0(y_i)/Z_z(\xi_i)\) for some point on the segment joining the two roots, provided its derivative is bounded away from zero there. A coincident root then gives zero displacement. Writing \(d=d_2-d_1\) and squaring the two half-gaps yields exactly
\begin{equation}
 D_Z-D_Y=d(h+d/4). \label{eq:direct}
\end{equation}
Both formulas use the same small difference \(B_0\). The first needs the reference slopes at its two roots, and the second needs target slopes along the connecting segments. The symmetric contour theorem retains its meaning even where those derivative divisions cease to be available.

\section{An obstruction to monotonicity on either parity}
\label{sec:parity}

At a fixed nonreal zero the leading quadrature error has a rotating
complex phase. Its projection onto the imaginary part of the root
displacement determines the signed change of the counted discriminant
\(D\) \cite[\S 2.2, Eq.~(2.4)]{agc}. The compact-source construction
below uses this phase to test monotonicity in the truncation
degree \cite{jh05}.

\Needspace{15\baselineskip}
\begin{proposition}[Displacement for a compact source]
\label{p:displacement}
Let \(\mu\) be an even probability measure with infinite support in
\([-K,K]\), and take \(F(z)=\int e^{izu}\,d\mu(u)\) and its
\(N\)-site Jacobi readout \(G_N\). Suppose a fixed circle centered on
the real axis encloses exactly two zeros of \(F\), a simple pair
\(\zeta=x+iy\) and \(\bar\zeta\) with \(x,y>0\), and that its boundary contains no zero.
For all sufficiently large \(N\), it encloses a simple pair
\(\zeta_N,\bar\zeta_N\) of \(G_N\), with \(\zeta_N\to\zeta\)
and a common zero-free collar about the circle. Put
\[
 A_N=\frac{h_N^*|\zeta|^{2N}}{(2N)!},\qquad
 \theta=\arg\zeta,\qquad D_N=-(\operatorname{Im}\zeta_N)^2 .
\]
Then the displacement and the two-degree increment are
\begin{align}
 \zeta_N-\zeta
   &=\frac{(-1)^N h_N^*\zeta^{2N}}{(2N)!\,F'(\zeta)}
     +O(A_N/N), \label{p:root-shift}\\
 \frac{D_{N+2}-D_N}{A_N}
   &=2y(-1)^N\operatorname{Im}
      \left(\frac{e^{2iN\theta}}{F'(\zeta)}\right)+O(N^{-1}).
      \label{p:degree-sign}
\end{align}
The implied constants are uniform in \(N\) with the source and
circle fixed.
\end{proposition}

\begin{proof}
The moment Gram matrices are positive definite by infinite support.
Multiplication by \(u\) on \(L^2(\mu)\) has norm at most \(K\), so
its finite Jacobi compressions satisfy \(\|J_N\|\le K\) and the
off-diagonal links satisfy \(0<\beta_j\le K^2\). In particular
\(0<h_N^*=\prod_{j=1}^N\beta_j\le K^{2N}\).

All interpolation knots lie in \([-K,K]\), so the Gaussian quadrature
error representation \eqref{eq:full-phase} applies on every fixed
complex disk. Write
\[
 H_N(z):=F(z)-G_N(z)=L_N(z)E_N(z),\qquad
 L_N(z)=(-1)^N\frac{h_N^*z^{2N}}{(2N)!},
 \qquad E_N(z)=\mathbb E e^{izX_N}.
\]
Here \(X_N\) is a uniform simplex average of the doubled Gauss
nodes and an independent source point \(U\) of law
\(P_N(u)^2\,d\mu(u)/h_N^*\), where \(P_N\) is monic. Reflection
makes \(X_N\) symmetric and \(|X_N|\le K\).
For the \(m=2N+1\) simplex coordinates \(S_j\), the identity
\(\mathbb E(S_j S_k)=(1+\mathbf 1_{j=k})/[m(m+1)]\) gives
\[
 \mathbb E X_N^2
   =\frac{2(\mathbb E U^2+\operatorname{tr} J_N^2)}{(2N+1)(2N+2)}
   \le \frac{K^2}{2N+1},
\]
since the doubled nodes sum to zero and have total squared sum
\(2\operatorname{tr}J_N^2\) \cite[\S 7, second moment]{vcd}.
Taylor's integral remainder, with
\(\mathbb E X_N=0\), now yields
\begin{equation}
 |E_N(z)-1|
 \le \frac{K^2|z|^2 e^{K|\operatorname{Im} z|}}{2(2N+1)},\qquad
 \sup_{|z|\le R}|H_N(z)|
 \le U_N(R):=\frac{e^{KR}h_N^*R^{2N}}{(2N)!}.
 \label{p:compact-bound}
\end{equation}

Choose disjoint closed isolating disks around the two zeros inside
the circle, arranged in conjugate positions. Choose \(R\) to
cover the circle, a closed zero-free collar of \(F\) and a fixed
positive margin around the disks. Since \(U_N(R)\to0\), Rouch\'e's
theorem gives counts \(2,1,1\) on these three circles, and uniform
convergence gives the common cleared collar. Let \(d_N=\zeta_N-\zeta\).
The positive minimum of \(|F(z)/(z-\zeta)|\) on the upper disk
gives \(d_N=O(U_N(R))\), using
\(F(\zeta_N)=H_N(\zeta_N)\). The fixed margin and Cauchy's
estimate bound \(H_N'\) on the disk by a constant times \(U_N(R)\).
Taylor expansion of \(F\) and integration of \(H_N'\) along the
displacement therefore give
\[
 F'(\zeta)d_N=L_N(\zeta)E_N(\zeta)+O(U_N(R)^2).
\]
The absolute disk bound retains the squared norm \(h_N^*\).
After division by the root scale it gives
\[
 \frac{U_N(R)^2}{A_N}
   =\frac{e^{2KR}h_N^*}{(2N)!}
       \left(\frac{R^2}{|\zeta|}\right)^{2N}
   \le \frac{e^{2KR}}{(2N)!}
       \left(\frac{KR^2}{|\zeta|}\right)^{2N}
   =o(N^{-1}).
\]
Together with \eqref{p:compact-bound} this proves
\eqref{p:root-shift}. Expanding \(D_N=-(y+\operatorname{Im}d_N)^2\)
and using \(A_N=o(N^{-1})\) gives
\[
 D_N+y^2=-2y(-1)^N A_N\operatorname{Im}
   \left(\frac{e^{2iN\theta}}{F'(\zeta)}\right)+O(A_N/N).
\]
Furthermore,
\[
 \frac{A_{N+2}}{A_N}
  =\frac{\beta_{N+1}\beta_{N+2}|\zeta|^4}
   {(2N+1)(2N+2)(2N+3)(2N+4)}
  \le \frac{K^4|\zeta|^4}{(2N+1)(2N+2)(2N+3)(2N+4)} .
\]
Thus \(D_{N+2}+y^2=O(A_NN^{-4})\), which proves
\eqref{p:degree-sign} and fixes the sign of its leading term.
\end{proof}

\begin{theorem}[Both signs occur on both parities]
\label{p:oscillation}
Set \(T(u)=(1-|u|)_+\) and take the continuous even density
\begin{equation}
 d\mu(u)=\frac{T(u)+[T(u-1)+T(u+1)]/4}{3/2}\,du .
 \label{p:density}
\end{equation}
Let \(D_N\) be the discriminant of the Jacobi pair counted inside
\(\lvert z-\pi\rvert=3/2\) for all sufficiently large \(N\). Put
\(a=\operatorname{arcosh}2\) and
\(A_N=h_N^*(\pi^2+a^2)^N/(2N)!\). For each \(p\in\{0,1\}\),
\begin{align}
 \limsup_{\substack{N\to\infty\\N\equiv p\pmod 2}}
     \frac{D_{N+2}-D_N}{A_N}
   &\ge \frac{a(\pi^2+a^2)}{2\sqrt{3}}, \notag\\
 \liminf_{\substack{N\to\infty\\N\equiv p\pmod 2}}
     \frac{D_{N+2}-D_N}{A_N}
   &\le -\frac{a(\pi^2+a^2)}{2\sqrt{3}} .
 \label{p:limbounds}
\end{align}
Consequently each parity contains infinitely many strict increases
and infinitely many strict decreases of \(D_N\).
\end{theorem}

\begin{proof}
Each tent has integral one. The density is positive on \((-2,2)\)
and has support \([-2,2]\). Integrating the central tent and
translating the two others gives
\[
 F(z)=\left[\frac{\sin(z/2)}{z/2}\right]^2
       \frac{1+\cos(z)/2}{3/2}.
\]
The squared sinc has zeros at the nonzero multiples of \(2\pi\).
The other factor has simple zeros at \((2k+1)\pi\pm ia\),
\(k\in\mathbb Z\). Since
\(\cosh(3/2)>1+(3/2)^2/2>2\), the two zeros
\(\zeta=\pi+ia\) and \(\bar\zeta\) lie strictly inside the
circle. Each other zero has horizontal distance at least \(\pi\)
from its center. In particular Proposition~\ref{p:displacement}
applies and \(0<4\theta=4\arctan(a/\pi)<\pi\).

At \(\zeta\), \(\sin\zeta=-i\sqrt{3}\) and
\(\sin^2(\zeta/2)=\cosh^2(a/2)=3/2\), giving
\(F'(\zeta)=2i\sqrt{3}/\zeta^2\). Formula \eqref{p:degree-sign}
becomes
\[
 \frac{D_{N+2}-D_N}{A_N}
   =-(-1)^N\frac{a(\pi^2+a^2)}{\sqrt{3}}
      \cos\bigl(2(N+1)\theta\bigr)+O(N^{-1}).
\]
On \(N=p+2j\) the normalized leading term divided by
\(a(\pi^2+a^2)/\sqrt{3}\) is a sinusoid \(b_j\) with angular step
\(4\theta\). Its first two Ces\`aro moments are \(0\) and \(1/2\).
Indeed the finite geometric sums at frequencies \(4\theta\) and
\(8\theta\) are bounded because neither frequency is a multiple of
\(2\pi\). For \(B=\limsup b_j\) and \(\epsilon>0\), eventually
\((b_j+1)(B+\epsilon-b_j)\ge0\), so
\[
 b_j^2\le B+\epsilon+(B+\epsilon-1)b_j .
\]
Taking Ces\`aro means and letting \(\epsilon\) decrease to zero
gives \(B\ge1/2\). The same argument for \(-b_j\) gives
\(\liminf b_j\le-1/2\). This proves \eqref{p:limbounds}
for both parities, with \(A_N>0\).
\end{proof}

The example excludes an eventual monotonicity inference based on
positivity, evenness, the cosine observation and a fixed count-two
chart for continuous compactly supported sources at time zero. The
canonical problem retains the full theta profile and its square
arithmetic \cite[\S 1.2, Eq.~(1.1)]{agc}.

\section{Uniform control over physical time}\label{sec:time}
To apply Theorem~\ref{th:signed} on a closed time interval, one needs a joint approximation to \(F-G_N\). The two fields move by amounts which can be much larger than their difference. Their Taylor coefficients may be subtracted before taking the range over time. We give the analytic justification for this procedure, since the dependence of the last retained link on the omitted links is essential.

\subsection{The complete field}
Anchor at \(t=0\), put \(t=Ts\) with \(0\le s\le1\), and write \(w=z-a\). With \(c_m=[w^m]F(0,a+w)\), let \(U(t,z)=H(t,z)/H(0,0)\) and \(g(t)=U(t,0)\). By the heat equation the coefficients of \(U\) and its mass are
\begin{align}
 [s^kw^m]U(Ts,a+w)&=\frac{(-T)^k(m+2k)!}{k!\,m!}\,c_{m+2k}, \notag\\
 [s^k]g(Ts)&=T^km_{2k}(0)/k! . \label{eq:taylor-source}
\end{align}
Since \(g(0)=1\), formal division determines the coefficients of \(F=U/g\). The xi identity in \eqref{eq:source} and the text following it supplies both sets of coefficients at the anchor.

Here is one way to justify truncation. Fix \(\rho>T\), a source disk \(|w|\le R_0\) and a working disk \(|w|\le R<R_0\). A sufficient condition for the chosen \(\rho\) is a strict upper bound below one for \(|g(t)-1|\) on \(|t|\le\rho\). It can be tested by integrating \((e^{\rho u^2}-1)\,d\mu_0\), and the numerator can be bounded by integrating \(e^{\rho u^2}\cosh(R_0u)\,d\mu_0\). A lower bound for \(|g|\) then gives a common upper bound \(M_F\) for \(|F|\) on the polydisk. Both integrals can be bounded without dropping a theta mode. For \(u\ge0\), \(n^4\le16^{n-1}\) and \(n^2-1\ge3(n-1)\) give
\[
 \Phi(u)\le \frac{2\pi^2}{1-16e^{-3\pi}}\exp(9u-\pi e^{4u}).
\]
Below \(u=1\) the untilted cosh integral is bounded by the full value \(F(0,iR_0)\), and the tilt factor costs at most \(e^\rho\). Above that point the displayed envelope dominates the tilted integrand explicitly. With \(\theta=T/\rho\), \(\lambda=R/R_0\), retention of \(0\le k\le K\) and \(0\le m<L\) gives a total remainder at most
\begin{equation}
 M_F\left\{\frac{\theta^{K+1}}{1-\theta}
        +\frac{\lambda^L}{(1-\lambda)(1-\theta)}\right\}.
 \label{eq:source-tail}
\end{equation}
The first term bounds the whole time tail by Cauchy's estimate at each fixed spatial point. The second sums the spatial tail of the retained time coefficients by Cauchy's estimate in both variables.

\subsection{Analytic continuation of the moving head}
The complex time remainder of \(G_N\) requires a disk on which the \emph{true} Jacobi coefficients are analytic. Define a tilted Gram matrix in the orthonormal basis at the anchor,
\[
 C_{ij}(t)=\int p_i(u;0)p_j(u;0)e^{tu^2}\,d\mu_0(u),\qquad
             0\le i,j<N .
\]
Thus \(C(0)=I\), and the entries are analytic in \(t\). The polynomials in both factors are fixed and real. This bilinear integral is its analytic continuation when \(t\) leaves the real axis.

\begin{proposition}[A Gram bound for the moving head]\label{prop:complex-head}
Suppose \(\|C(t)-I\|_2\le\epsilon<1/2\) throughout \(|t|\le\rho\).
Let \(L_0\) bound the maximum absolute row sum of \(J_N(0)\), let \(\sigma=\epsilon/(1-\epsilon)\), and set \(d=L_0\sigma/(1-\sigma)\). The physical links extend analytically to the disk with their positive values at zero, and
\begin{align}
 \|J_N(t)-J_N(0)\|_2&\le d,\notag\\
 |G_N(t,z)|&\le \exp\{|\Im z|L_0+|z|d\}. \label{eq:duh}
\end{align}
\end{proposition}
\begin{proof}
Each leading principal block obeys the same perturbation bound. For its last Schur pivot \(S_j\), the diagonal perturbation is bounded by \(\epsilon\), each off-diagonal block by \(\epsilon\), and the inverse of the preceding leading block by \(1/(1-\epsilon)\). Thus \(|S_j-1|\le\epsilon+\epsilon^2/(1-\epsilon)=\sigma<1\). The pivot equals the ratio of the tilted unnormalized monic norm of degree \(j\) to its time-zero value. Indeed the Schur complement subtracts the projection on lower degrees, and a triangular change of basis preserves the leading coefficient. The scalar mass normalization cancels when consecutive norms are divided. Hence
\[
 \frac{a_j(t)}{a_j(0)}
  =\exp\!\left(\tfrac12(\log S_j-\log S_{j-1})\right),\qquad 1\le j<N .
\]
The logarithms are analytic around \(1\) and have modulus at most \(-\log(1-\sigma)\). The change in this ratio has modulus at most \(\sigma/(1-\sigma)\). The absolute row and column sums of the matrix change are each bounded by \(d\), which proves the first inequality. For \(0\le v\le1\), expand \(e^{ivzJ_N(t)}\) by variation of constants about \(J_N(0)\). Along the reference evolution its norm is bounded by \(e^{v|\Im z|L_0}\) by self-adjointness at the anchor. Grönwall's inequality gives~\eqref{eq:duh} at \(v=1\).
\end{proof}

The hypothesis itself can be acquired from the complete source. On \(|u|\le1\), Cauchy--Schwarz and orthonormality bound the operator change by \(e^\rho-1\). On the complement bound it by the integral of \(e^{\rho u^2}\sum_{j<N}|p_j(u;0)|^2\) against \(d\mu_0\), using \(|e^{tu^2}-1|\le e^{\rho u^2}-1\). The roots of \(P_j(\cdot;0)\) are the eigenvalues of a leading compression and have magnitude at most \(L_0\). For \(|u|\ge1\) this gives
\[
 \sum_{j<N}|p_j(u;0)|^2\le
 \frac{N(1+L_0)^{2N-2}}{\min_{j<N}h_j(0)}\,|u|^{2N-2}.
\]
The theta envelope controls the remaining integral. This use of source eigenvalues imposes no ordering of Fourier zeros.

Finally, iteration of~\eqref{eq:volterra} shows by induction that order \(k\) of \(\beta_j(Ts)\) needs initial links only through \(\beta_{j+k}(0)\): each differentiation reaches at most one neighbor farther along the chain. The order-\(K\) head therefore uses anchor moments through \(2(N-1+K)\). Taylor coefficients in \(s\) of the first \(P\) spatial exponential terms follow by multiplying tridiagonal matrices, or by summing paths, modulo \(s^{K+1}\). For \(|z-a|\le R\) put \(Z_0=|a|+R\), \(v=Z_0(L_0+d)\). The absolute joint remainder is at most
\begin{equation}
 e^{RL_0+Z_0d}\frac{\theta^{K+1}}{1-\theta}
       +\frac{e^{v}v^P}{P!(1-\theta)} .
 \label{eq:head-tail}
\end{equation}
The first term is Cauchy's time estimate for the full matrix readout. The exponential tail is bounded by \(e^vv^P/P!\) throughout the complex time disk. Projecting that tail onto the retained time orders costs at most \(1/(1-\theta)\), giving the second term. Subtracting the two coefficient arrays for \(F\) and \(G_N\) at a common \(s,w\) preserves their cancellations. Their complete remainder bounds add. A radius margin of one between these working disks and a smaller disk also bounds the first spatial derivative of each remainder there by Cauchy's formula.

\section{Supporting evidence on one canonical circle}\label{sec:evidence}
Take \(Y=G_{24}\), \(Z=F\), \(\Gamma=\{|z-36|=9\}\), and \(0\le t\le T=10^{-4}\). The measured quantity is
\(\Delta(t)=D_F(t)-D_{24}(t)\). The earlier complete-theta certificate \cite{hdc} counts two zeros in the rectangle \((26,46)+i(-2,2)\) at time zero. Since the time-zero xi zeros lie in \(|\Im z|\le1\), every possible zero in either disk of radius \(9\pm1/64\) is in that rectangle. Two real sign changes exhaust its count. The following additional bounds are needed to extend this geometry to the slab and determine the sign of the correction \cite{jh07,jh08}.

The construction in Section~\ref{sec:time} used \(\rho=1/20\), \(R_0=24\), \(R=10\) and time degree \(K=10\). Initial moments through order 66 supply the moving head through that order. Its Gram perturbation is below \(0.051272\). With 144 centered spatial terms for the source and 192 origin exponential terms for the head, the respective joint tails on radius ten are below \(4.750\times10^{-29}\) and \(8.876\times10^{-27}\). Outward coefficient intervals include all rounding. The polynomial for \(F-G_{24}\) is formed before evaluation over the time interval.

On both collar boundaries a Taylor enclosure covers each of 512 angular sectors. Nonvanishing for every \(t\) preserves the source counts by homotopy, and the strict relative bound transfers both counts to the head. Sign changes in \((28,29)\) and \((42,43)\) for each field at every \(t\) place both simple real roots inside radius eight. Equality of the two circle counts excludes zeros from the annulus. Maximum modulus applied to the quotient and reciprocals therefore gives the uniform bounds
\[
\begin{gathered}
  |F|>1.4331937907\times10^{-5},\qquad
  |G_{24}|>1.4331935452\times10^{-5},\\
  |(F-G_{24})/G_{24}|<1.715356\times10^{-7}.
\end{gathered}
\]
In particular \(b=1\). The exact choice \(q_0=q\) gives zero modeling error \(e\). Its numerical evaluation errors are included in the outward intervals. The two terms of the nonlinear allowance in Theorem~\ref{th:signed} sum to less than \(2.846816\times10^{-12}\).

\begin{table}[htb]
\centering\small
\caption{Certified enclosures of the same correction over the entire interval. Widths use the stored endpoints before outward decimal rounding.}\label{tab:certificate}
\begin{tabular}{@{}lcc@{}}
\toprule
Method & Interval for \(10^8\Delta(t)\) & Bound on width of \(\Delta\)\\
\midrule
Signed contour & \((-4.958280826,-4.956734690)\) & \(1.546135\times10^{-11}\)\\
Direct displacement & \((-4.957996104,-4.957019318)\) & \(9.767860\times10^{-12}\)\\
\bottomrule
\end{tabular}
\end{table}

Table~\ref{tab:certificate} uses \eqref{eq:residue} and \eqref{eq:direct} with the same source, head, root intervals and numerator. In particular, the direct test retains the correlation between the two fields at each time. The slope bounds on both connecting segments exclude zero. The unsigned logarithm bound allows \(|\Delta|<2.624495\times10^{-5}\), leaving the sign unresolved, although it already proves positivity of the absolute gap on this known chart. The signed and direct evaluations show that the head overestimates the squared half-gap throughout the interval. The direct enclosure lies strictly inside the signed one.

The time-zero calculation alone encloses \(\Delta(0)\) near \(-4.957710484497283\times10^{-8}\) with direct width below \(2.073\times10^{-35}\). Subtracting this saved interval from the same-time direct evaluation at \(T\) gives
\[
 4.05561835\times10^{-12}<\Delta(T)-\Delta(0)<4.05588776\times10^{-12}.
\]
This measures a decrease in the magnitude of the error between the two endpoints. Its derivative in the interior has no certified sign here. The full-interval enclosure came from the range of a joint analytic model. The repository records the outward endpoints, evaluator versions, counts, and the matched internal replay \cite{jh07,jh08}.

A simple event check distinguishes the two possible sign decisions. Take
\(Y=(z-36)^2-2t-1\) and \(Z=(z-36)^2-2t\). Both obey the unnormalized heat equation, and their squared half-gaps are \(2t+1\) and \(2t\). The correction is exactly \(-1\), even at the target double at \(t=0\). The strict absolute-gap condition following Theorem~\ref{th:signed} fails there, as do distinct target brackets and a target slope floor. The contour coordinates remain defined. This example isolates the distinction between an approximation error and a collision margin.

\section{Outstanding challenges}\label{sec:open}
\paragraph{How much of the source is needed for the sign?}
The finite degree formula suggests a first concrete experiment in information reduction. Keep the same circle, time interval and target \(D_F-D_{24}\), but form \(q_0=(G_{26}-G_{24})/G_{24}\) from finite readouts. To predict the sign one must independently bound the rest, \(e\) in \eqref{eq:remaining}, and bound the complete quotient by \(\eta<1\) on the collar. Here \(r(2r-b)=153\). If \(J_0=-\langle(z-c_{24})q_0\rangle\) and the nonlinear allowance stay near the observed values, an error appreciably below \(4.96\times10^{-8}/153\), about \(3.2\times10^{-10}\) in relative units, would suffice. The actual margin must be computed for \(q_0\). Using the complete numerator of Section~\ref{sec:evidence} to determine that margin would spend the information the finite step is intended to save. A meaningful comparison gives a direct evaluator the same finite readouts, real bounds, and residual allowance.

\paragraph{Can the degree estimates be iterated?}
A single signed correction supplies neither a common count at arbitrarily large \(N\) nor a rule for its sign at the next degree. One possible sufficient mechanism is to control \(u_N=\max(0,-D_N)\) on a fixed parity by
\[
 u_{N+2}\le (1-\alpha_N)u_N+\varepsilon_N,\qquad
 0<\alpha_N\le1,\qquad \sum_N\alpha_N=\infty,\qquad
 \varepsilon_N/\alpha_N\longrightarrow0 .
\]
The sums and limits are on that parity and \(\varepsilon_N\ge0\). Indeed, for any \(\delta>0\), eventually \(\varepsilon_N\le\delta\alpha_N\), so \((u_{N+2}-\delta)_+\le(1-\alpha_N)(u_N-\delta)_+\). The product tends to zero and then \(\delta\downarrow0\) gives \(u_N\to0\). To infer \(D_F\ge0\), one must also identify the same limiting cluster and show \(D_N\to D_F\) on its counted chart. A source-dependent inequality of this kind is open. Theorem~\ref{p:oscillation} rules out a general monotonicity shortcut for compact positive sources. A strictly positive lower bound would demand more than vanishing negative part.

\paragraph{Can time transport give new coverage?}
The full forcing in~\eqref{d:weighted-gap} has a signed endpoint representation on every subinterval. Useful backward continuation requires a common two-zero chart, an initial absolute margin, the true exterior weight, and a bound on the worst cumulative loss before a possible event, as in \cite[(4.3)]{agc}. The source must control these jointly if cancellations are to survive. A change of chart needs a verified identification of the enclosed pair, and a cluster of larger multiplicity needs more coordinates. The uniform finite-interval calculation here quantifies approximation error in a region with a known prior count. Extending this information to other heights or longer times, or using it to improve a bound for the de Bruijn--Newman constant, remains open.

\medskip
\noindent\textit{Acknowledgment and status.} This working draft was prepared for Edward Baker with OpenAI Codex assistance in derivation, exposition and internal checking. Bounded drafting and checking contributors were dispatched as gpt-6.2-sol with ultra reasoning effort. The public session record does not identify the lead configuration. The assessments share the repository's source assumptions and constitute internal review.

\begin{thebibliography}{99}
\footnotesize
\setlength{\itemsep}{1.7pt plus 0.5pt}
\bibitem{agc} E.~Baker, \emph{Arithmetic residuals and backward gap continuation in the Newman heat flow}, parent investigation manuscript, 19 November 2026. Equation numbers cited in the text refer to this version.
\bibitem{vcd} E.~Baker, \emph{Jacobi dynamics, signed error, and cluster transport for the Newman heat flow}, working manuscript, 17 November 2026. See the sections on Gaussian quadrature error, phase bounds and gap transport.
\bibitem{cm} E.~Baker, \emph{Exterior moment hierarchies and coupled gap dynamics in the Newman heat flow}, working manuscript, 18 November 2026.
\bibitem{gt} E.~Baker, \emph{One canonical signed gap estimate}, cluster gap transport investigation, synthesis note 16, 16 November 2026.
\bibitem{jh02} E.~Baker, \emph{The forced integer hierarchy for the moving Jacobi head}, Jacobi--integer hierarchy investigation, note 02, 19 November 2026.
\bibitem{jh03} E.~Baker, \emph{Signed degree increment and physical-time compatibility}, same investigation, note 03, 19 November 2026.
\bibitem{jh04} E.~Baker, \emph{Signed gap correction on a counted chart}, same investigation, note 04, 19 November 2026.
\bibitem{jh05} E.~Baker, \emph{A compact-source control for parity monotonicity}, same investigation, note 05, 19 November 2026.
\bibitem{jh07} E.~Baker, \emph{Measuring the Jacobi gap correction on the time-zero circle}, same investigation, note 07, numerical record 05 and assessment 03, 19 November 2026.
\bibitem{jh08} E.~Baker, \emph{The degree-24 gap correction on a physical-time interval}, same investigation, note 08, numerical record 06 and assessment 04, 19 November 2026.
\bibitem{hdc} E.~Baker, \emph{A complete theta wall and coefficient certificate in the frozen finite box}, height-dependent heat contours investigation, note 13 and numerical record 07, 23 October 2026.
\bibitem{polymath} D.~H.~J. Polymath, \emph{Effective approximation of heat flow evolution of the Riemann \(\xi\) function, and a new upper bound for the de Bruijn--Newman constant}, 2019. \href{https://arxiv.org/abs/1904.12438v2}{arXiv:1904.12438v2}.
\bibitem{vanassche} W.~Van Assche, \emph{Orthogonal polynomials, Toda lattices and Painlev\'e equations}, Physica D \textbf{434} (2022), 133214, \S2.3. \href{https://arxiv.org/abs/2202.11017v1}{arXiv:2202.11017v1}.
\bibitem{dlmf} NIST Digital Library of Mathematical Functions, \href{https://dlmf.nist.gov/3.5}{\S3.5(v)--(vi)}, Gauss quadrature and its eigenvalue characterization.
\end{thebibliography}

\end{document}
